% Generated draft: requires mathematical and source review.
\documentclass[11pt]{article}
\usepackage[margin=25mm]{geometry}
\usepackage{fontspec}
\setmainfont{Latin Modern Roman}
\usepackage{amsmath,amssymb,mathtools,unicode-math}
\setmathfont{Latin Modern Math}
\usepackage{enumitem,needspace,xurl,hyperref}
\setlength{\parindent}{0pt}
\setlength{\parskip}{5pt}
\setlength{\emergencystretch}{3em}
\newcommand{\N}{\mathbb{N}}
\newcommand{\Z}{\mathbb{Z}}
\newcommand{\Q}{\mathbb{Q}}
\newcommand{\R}{\mathbb{R}}
\newcommand{\C}{\mathbb{C}}
\newcommand{\F}{\mathbb{F}}
\newcommand{\A}{\mathbb{A}}
\newcommand{\PP}{\mathbb{P}}
\newcommand{\E}{\mathbb{E}}
\newcommand{\eps}{\varepsilon}
\newcommand{\dd}{\,\mathrm{d}}
\begin{document}
\section*{UnsolvedMath ID: 30002095}
\textbf{Source problem code:} catalogue entry
\subsection{Definitions and mathematical statement}
Let \(x\) denote the signal mentioned in the supplied record. The record does not specify its domain, codomain, function space, or admissible signal class. The terms \emph{low-pass measurements}, \emph{noise}, \emph{cutoff resolution}, \emph{finer features}, and \emph{recovery accuracy} are not given mathematical definitions.

\textbf{Source question.} Given noisy low-pass measurements of \(x\), determine how accurately its finer features can be recovered by super-resolution as the desired resolution becomes finer than the cutoff resolution.

\textbf{Unresolved formulation.} This is a qualitative research question, not a fully specified mathematical assertion. A precise formulation requires a measurement operator and observation domain, an admissible signal class, a noise model and noise magnitude, a definition of the features to be recovered, a resolution parameter and its relation to the cutoff, and an error criterion. It must also specify whether accuracy concerns a particular reconstruction method or optimal performance over a class of methods, together with the relevant quantifiers. None of these choices is fixed by the supplied record. Consequently, no definite recovery bound, asymptotic rate, or quantified theorem can be stated without adding hypotheses.
\subsection{Short English statement}
The question asks how recovery accuracy changes when super-resolution seeks details finer than the measurement cutoff. The supplied record leaves the signal model, noise, target features, and accuracy criterion unspecified, so it does not determine a precise mathematical claim. Its supplied \emph{open} label is a curation status, not an independently verified mathematical conclusion.
\subsection{Sources}
\begin{itemize}
\item Local UnsolvedMath registry: supplied complete source record, permanent ID \(30002095\), source problem code \texttt{catalogue entry}, title \emph{Super-Resolution Accuracy beyond the Cutoff Scale}, in the Oberwolfach Reports Open Problems collection. The project identifies \texttt{data/unsolvedmath/problems.json} and \texttt{catalog/problems.json} as local registry files; no independent inspection is asserted here.
\item Supplied report citation: \emph{Applied Harmonic Analysis and Sparse Approximation (2013)}, Oberwolfach Reports, DOI \texttt{10.4171/owr/2012/29}, pp.\ 1759--1843. Supplied source URL and literature-source link, labeled ``original OWR report'': \url{https://doi.org/10.4171/owr/2012/29}. The background attributes the question to workshop open-problem discussions documented in Oberwolfach Reports; no individual proposer is supplied.
\item Literature assessment dated \texttt{2026-08-21}: the supplied curation labels the item open and states that the context-free question does not specify a model, observations, or target theorem, so no determinate solved or disproved claim can be identified. This is supplied dated curation, not independently verified literature assessment; the supplied background requests independent verification of time-sensitive status and cited claims.
\item The registry labels the statement \texttt{corrected\_verified} and reports manual semantic, typographic, and LaTeX review of the cleanup. This supplied verification claim does not resolve the missing mathematical specifications.
\end{itemize}
\end{document}
